% discussion.tex
%
% This is the ONLY section permitted to interpret. Items deliberately held
% back from results.tex and discharged here:
%   - condition number as a scale-sensitive quantity (supported by this
%     study's own normalization-operator measurement);
%   - the ceiling observation as the organizing account of why the four
%     results cohere.
%
% CLAIM CEILING (governance rule, ../../27_One_Sentence_Paper.md). These are
% CUT, not softened, wherever they appear: "the mechanism does not transfer",
% "geometry stops explaining reliability", "foundation models invalidate the
% account", "the phenomenon is absent", "fails", "does not generalize".
%
% Distinction that must be preserved throughout (PI correction, 2026-08-06):
% claims are about what OUR EXPERIMENTS SUPPORT, never about what exists in
% the world. Write "our experiments provide limited support for testing the
% account in this regime", never "there is little for the account to explain".
%
% VERIFICATION DEBT: \cite{tran2026decodable} carries the comparison figure
% on which the two-orders-of-magnitude statement rests.

\section{Discussion}

\subsection{Principal Findings}

We tested whether an association reported across representations---between
in-distribution geometry and distance-based estimator performance---can
account for a previously reported phenomenon, the probe--reliability
dissociation, when geometry is intervened upon within a single fixed frozen
backbone. Under the evaluated backbone, layer, and task, that explanatory
reading received limited empirical support. The previously reported
dissociation was reproduced in direction with high statistical confidence,
but at a magnitude approximately two orders of magnitude smaller than
originally reported. Neither preregistered geometry metric showed a
Holm-corrected association with estimator performance, and a preregistered
comparison that matched spectral-decay slope across differing nominal
intervention strengths did not reach corrected significance in either
direction. Three of sixteen intervention-type-by-geometry interactions
remained Holm-significant after conditioning on geometry which, under the
preregistered criterion, is inconsistent with full mediation.

\subsection{A Coherent Reading of the Four Analyses}

The four analyses address different aspects of the question, and they are
consistent with a single observation: on this backbone, distance-based
estimators were already operating near ceiling before any intervention, with
Reference-arm AUROC of $0.9984$ for Mahalanobis and $0.9979$ for
nearest-neighbour. The domain shift studied here was, for these estimators,
already highly separable.

This observation makes the remaining results easier to situate. Where an
estimator has almost no headroom, there is correspondingly little variation
in its performance for a geometry metric to covary with, which is consistent
with the absence of corrected-significant associations for LID and
spectral-decay slope. It is likewise consistent with the causal comparison,
in which Mahalanobis and nearest-neighbour showed no separation between any
configuration, and only the two estimators retaining measurable headroom---
cosine and kernel density---separated at all. That two of the four estimators
were insensitive to every configuration is itself a finding rather than an
absence of one: it establishes that on this representation the two strongest
distance-based detectors are unaffected by interventions of the magnitude
studied here, which is directly relevant to whether geometry-targeted
correction is warranted in deployment. And in the interaction models,
where residual variation is at the fourth decimal place, a statistically
detectable interaction term need not correspond to a substantively meaningful
effect; the three Holm-significant coefficients were of order $10^{-4}$ to
$10^{-2}$.

We emphasize what this reading does and does not assert. It is a statement
about the sensitivity our experiments could achieve, not about the account
itself. An evaluation conducted where the phenomenon of interest is small
provides limited power to test an explanation for that phenomenon; the
appropriate conclusion is that this study offers limited empirical support
for the account in this regime, not that the account is incorrect, nor that
the geometric quantities it names are unimportant in regimes where the
phenomenon is large.

\subsection{Condition Number}

Condition number was the one geometric quantity associated with all four
estimators at corrected significance, while the two preregistered metrics
were not. This study provides direct evidence bearing on how that association
should be read. When we measured the normalization-based operator, it changed
condition number by approximately $2.7\times10^{5}$ while leaving LID and
spectral-decay slope unchanged to within $10^{-3}$ and $10^{-6}$
respectively. Condition number therefore responds strongly to a
transformation that leaves both shape metrics untouched, consistent with it
behaving substantially as a scale-sensitive descriptor rather than as a
measure of the spectral shape at issue. This reading is reinforced by the
operator's own specification, which adjusts feature radii while holding
directions fixed~\cite{yu2510dissecting}. Its association with estimator
performance should accordingly not be read as support for the explanatory
hypothesis, and we do not treat it as such.

\subsection{Limitations}

Several limitations bound these findings.

The correlational result depends on how the pooled configurations are treated.
Holm-corrected significance for condition number holds when the $30$
configurations are treated as the unit of analysis; under an exact permutation
test over whole seeds none of the four estimator associations survives
correction. That sensitivity analysis is reported with the results and was not
preregistered. The two analyses are not in conflict so much as differently
powered: the seed-level test has only $5! = 120$ arrangements available and
therefore very little resolution, so it cannot adjudicate an association of this
size in either direction. What can be said is that the strength of the
correlational claim is contingent on the resampling unit, and readers who take
the seed as that unit should treat the association as suggestive rather than
established.

The evaluation covers one backbone, one representation layer, one diagnostic
task, and one domain-shift pair. That pair differs in imaging modality, not
only in acquisition conditions: ISIC 2019 comprises dermoscopic images and
PAD-UFES-20 close-up clinical photographs. A modality difference is a
stronger form of distribution shift than the device- or site-level shifts
that motivate distance-based reliability estimation in deployment, and it is
the most plausible reason the two populations are as separable as we observe.
Whether the near-ceiling estimator performance reflects a property of frozen
foundation-model representations in general, or the modality gap between
these particular populations, cannot be determined from this study. A
same-modality shift---for instance between dermoscopic cohorts acquired at
different sites or with different devices---would produce lower baseline
separability and a correspondingly more sensitive test. We regard
constructing such a comparison as a precondition for a stronger evaluation of
the explanatory hypothesis in this regime rather than as a refinement of the
present one.

The causal comparison was matched on spectral-decay slope alone. At the
calibrated exponent, LID remained $0.60$ below the target checkpoint's value,
so the comparison does not establish matching on both preregistered metrics
jointly. Because the account names both quantities, this is a genuine
restriction on what the comparison can support.

The manipulation operator's effect on LID was reproducible but
non-monotonic, rising at intermediate exponents before falling at full
strength, with the same profile in all five seeds. A monotone dose--response
was not assumed by the preregistered criterion, and the observed profile
satisfies it; but it means the operator does not provide a simple
one-directional LID control, which constrains the design of future
comparisons.

Finally, five seeds with an eight-comparison correction family afford limited
statistical power for the causal comparison, and we report it as
inconclusive at the significance-test level rather than as evidence in either
direction.

\subsection{Implications}

For practitioners deploying distance-based reliability estimation on frozen
foundation-model representations, these results provide no evidence that
geometry-targeted correction is required before such estimators can be used
in the configuration evaluated here. This is a statement about the evaluated
setting, and does not extend to configurations where baseline separability is
lower.

Methodologically, the value of this study lies less in any individual
statistic than in the sequence of checks that constrained them. Two elements
proved decisive. The manipulation check identified, before any causal result
was interpreted, that the operator originally specified for the comparison
did not move the quantities at issue---a finding that would otherwise have
surfaced only as an unexplained null. And preregistering the correction
family before execution meant that several uncorrected $p$-values below
$0.05$ could be reported as what they are, rather than selected after the
fact.

Rather than confirming an explanatory reading of the reported association,
our results delineate the boundary under which such a reading can currently
be supported: intervening on representation geometry within a single frozen
backbone yielded limited empirical support for it, under the evaluated
backbone, layer, and task.
